% 05_formalism.tex
% Generated by theory-scaffold for idea-slug: [IDEA_SLUG]
% This file is the authoritative formalism. Stage 6 (numerical-prototype) reads it
% directly and the reconciliation script checks symbol correspondence.

\documentclass[11pt]{article}
\usepackage{amsmath,amssymb,amsthm}
\usepackage{physics}
\usepackage{hyperref}
\usepackage{cleveref}
\usepackage[margin=1in]{geometry}

\title{Formalism for [IDEA\_SLUG]}
\author{Generated by theory-scaffold}
\date{\today}

\begin{document}
\maketitle

\section{Framework and assumptions}
\label{sec:framework}

The framework is [governing framework name, e.g., generalized Langevin / Fokker-Planck / variational].

\paragraph{Conservation laws.}
\begin{itemize}
    \item [list, e.g., total particle number conserved]
\end{itemize}

\paragraph{Symmetries.}
\begin{itemize}
    \item [list]
\end{itemize}

\paragraph{Assumptions.} Each assumption is labeled for cross-reference.
\begin{description}
    \item[(A1)] [assumption statement]. Valid in [regime]. Cost of violation: [qualitative / quantitative / catastrophic].
    \item[(A2)] ...
\end{description}

\section{Core result}
\label{sec:core}

The headline equation is
\begin{equation}
    \label{eq:core}
    [core equation]
\end{equation}
where [define every symbol on first use].

Physical interpretation: [1-2 sentences].

\section{Derivation}
\label{sec:derivation}

% Each step follows the template in stepwise_derivation_protocol.md.
% One conceptual operation per step. Verification pass per step.

\subsection{Step 1: [operation name]}

\paragraph{Starting point.}
\begin{equation}
    \label{eq:step1-in}
    [equation in]
\end{equation}

\paragraph{Operation.} [one sentence]

\paragraph{Result.}
\begin{equation}
    \label{eq:step1-out}
    [equation out]
\end{equation}

\paragraph{Verification.} [dimensional / limit / sanity]

\paragraph{Assumptions invoked.} (A1)

% Continue with Step 2, Step 3, ...

\section{Limits}
\label{sec:limits}

\subsection{Limit 1: [named limit, e.g., dilute limit $\phi \to 0$]}

In this limit, [describe what the theory should reduce to]. Taking [specific parameter limit] in \Cref{eq:core} yields [result], which matches [known reference], citing [reference]. \checkmark

\subsection{Limit 2: [named limit]}

[analogous]

\section{Dimensionless groups}
\label{sec:pi}

Applying the Buckingham Pi theorem to the variables [list], we identify [n] independent dimensionless groups:
\begin{align}
    \Pi_1 &= [definition] && [\text{physical meaning}] \\
    \Pi_2 &= [definition] && [\text{physical meaning}] \\
    &\vdots
\end{align}

The governing equation \Cref{eq:core} in dimensionless form is
\begin{equation}
    \label{eq:core-dimensionless}
    [dimensionless version]
\end{equation}

\section{Gray-box boundary (if applicable)}
\label{sec:graybox}

See \texttt{templates/graybox\_boundary.md} for the template. Specify:
\begin{itemize}
    \item Physics-prescribed terms: [list with references to Eq.~\ref{eq:core}]
    \item Learned terms: [list with per-component validation plans]
    \item Training data and ranges: [specify]
\end{itemize}

\section{Open questions for Stage 6}
\label{sec:open}

\begin{itemize}
    \item [question the numerical prototype must resolve]
\end{itemize}

\end{document}
